\documentclass[12pt]{article}
\usepackage[utf8]{inputenc}
\usepackage[T1]{fontenc}
\usepackage{hyperref}
\title{Pasi Language Resources for K-12 Curriculum Development}
\author{Compiled by Claude Sonnet 5.5 (claude-sonnet-5-5) for Andr\'as Kornai}
\date{}
\begin{document}
\maketitle

\section{Introduction}
Pasi (ISO 639-3 psq; Glottocode pasi1259) is a Sepik language of Papua New Guinea. The searches found only general reference entries, and no educational materials of any kind. Sources disagree on speaker numbers (360 in one, fewer than 154 in another, 700 for the people group in a third) and on location (Sandaun Province in linguistic sources, Gulf Province in Joshua Project), so the identification should be checked. Do not confuse it with the Adi variety also called Pasi.

\section{Primary Educational Resources}
None found.

\section{Linguistic and Reference Materials}
\begin{itemize}
\item Wikipedia, Pasi language: \url{https://en.wikipedia.org/wiki/Pasi_language}
\item Glottolog entry: \url{https://glottolog.org/resource/languoid/id/pasi1259}
\item Verbix endangered languages entry: \url{https://docs.verbix.com/EndangeredLanguages/Pasi}
\item ASJP word list entry: \url{https://asjp.clld.org/languages/PASI}
\end{itemize}

\section{Religious and Community Resources}
Joshua Project reports no Bible translation and no Christian resources in Pasi: \url{https://joshuaproject.net/languages/psq} and \url{https://joshuaproject.net/people_groups/19519}.

\section{Digital and Audio Resources}
None found.

\section{Curriculum Development Context}
Nothing indicates Pasi is taught in schools. Any curriculum would have to begin with documentation and community consultation.
\end{document}